\chapter[Fundamentação Teórica]{Fundamentação Teórica}
\label{cap:fundamentacao-teorica}

Os conceitos apresentados aqui cobrem o mínimo necessário para entender as escolhas metodológicas do Capítulo~\ref{cap:Metodologia}, sem incluir tutoriais gerais de aprendizado profundo.

\section{Formulação do problema}

Verificação de parentesco facial é uma tarefa de classificação binária. Dado um par de imagens de rosto, $x_a$ e $x_b$, o sistema decide se as duas pessoas têm ou não relação de parentesco. O modelo aprende uma função $f_\theta : \mathbb{R}^{H \times W \times 3} \to \mathbb{R}^d$ que mapeia cada face para um vetor de embedding de dimensão $d$. A similaridade entre dois embeddings é medida por cosseno, $s(x_a, x_b) = \cos(f_\theta(x_a), f_\theta(x_b))$, e a decisão final compara essa similaridade contra um limiar $\tau$. Pares com $s > \tau$ são classificados como parentes, os demais como não-parentes.

Os modelos supervisionados deste trabalho abordam a versão binária do problema. As linhas de base com VLMs, por outro lado, abordam uma variante multiclasse, em que a saída é uma das onze relações do FIW como pai-filho, mãe-filha, irmãos, avô-neto e demais. Essa diferença de tarefa importa para a comparação entre regimes e é discutida com mais detalhe no Capítulo~\ref{cap:Metodologia}.

\section{Redes siamesas e aprendizado por similaridade}

Uma rede siamesa é composta por dois ramos idênticos que compartilham pesos e processam, cada um, uma imagem do par. Os embeddings resultantes são comparados por uma medida de similaridade ou distância, e a rede aprende a aproximar pares positivos e afastar pares negativos no espaço de embeddings. Essa estrutura é central para verificação facial e foi adaptada com sucesso para parentesco no trabalho seminal de \citeonline{robinson2018visual}.

A medida de similaridade mais usada neste trabalho é a cosseno, definida sobre embeddings $L^2$-normalizados, que produz valores no intervalo $[-1, 1]$ e é invariante a magnitudes. A distância euclidiana é uma alternativa equivalente quando os embeddings têm norma constante. O conceito de margem aparece nas funções de perda discriminativas e quantifica a separação mínima exigida entre pares positivos e negativos no espaço de embeddings.

\section{Redes convolucionais e ConvNeXt}

Uma rede neural convolucional, ou CNN, aprende filtros locais que percorrem a imagem de entrada e produzem mapas de características. Camadas sucessivas combinam filtros de campo receptivo crescente, e operações de pooling reduzem a resolução espacial ao longo da hierarquia. As características das camadas iniciais capturam bordas e texturas, enquanto as camadas profundas capturam partes da face e estruturas semânticas. ResNet \cite{he2016deep}
é o bloco-base moderno em reconhecimento facial e fornece a coluna do Modelo 01 deste trabalho.

ConvNeXt \cite{liu2022convnet} revisita a CNN tradicional incorporando ideias do modelo Transformer, como particionamento da entrada em blocos, convoluções depthwise, blocos invertidos e normalização por camada em vez de batch. Os autores mostram que essa CNN modernizada atinge desempenho competitivo com Vision Transformers em ImageNet, mantendo a eficiência computacional das convoluções. No Modelo 03 deste trabalho, o ConvNeXt forma metade do duplo backbone e contribui com características locais e vieses indutivos convolucionais que complementam as representações globais do ramo Transformer.

\section{Vision Transformers e DINOv2}

O Vision Transformer \cite{dosovitskiy2020image} aplica o modelo Transformer originalmente proposto para texto \cite{vaswani2017attention}
diretamente sobre imagens. A imagem é dividida em patches não-sobrepostos, cada patch é projetado linearmente em um vetor, e os vetores são concatenados com embeddings de posição para formar a sequência de entrada do encoder Transformer. O bloco central do Transformer é a self-attention, definida por

\begin{equation}
\text{Attention}(Q, K, V) = \mathrm{softmax}\!\left( \frac{QK^\top}{\sqrt{d_k}} \right) V,
\label{eq:attention}
\end{equation}

onde $Q$, $K$ e $V$ são projeções lineares dos tokens de entrada chamadas respectivamente de query, key e value, e $d_k$ é a dimensão de cada cabeça de atenção. A atenção permite que cada token atue sobre todos os outros tokens da sequência, capturando dependências de longo alcance que CNNs precisam de muitas camadas para alcançar. No domínio facial, isso favorece a captura de relações entre partes distantes do rosto, como olhos e mandíbula, sem depender de campos receptivos hierárquicos.

DINOv2 \cite{oquab2024dinov2}
é um Vision Transformer treinado de forma auto-supervisionada sobre 142 milhões de imagens sem rótulos. O método combina destilação entre recortes aumentados de uma mesma imagem e um objetivo de máscara em nível de patch, e produz representações que se transferem bem para tarefas de granularidade fina sem ajuste fino completo. No Modelo 05 deste trabalho, o backbone DINOv2 fica congelado e a adaptação ao domínio é feita apenas pelos adaptadores LoRA descritos adiante.

Em contraste com backbones de uso geral, modelos especializados em reconhecimento facial usam arquiteturas e perdas voltadas à separação de identidades em larga escala. O AdaFace \cite{kim2022adaface}, baseado em uma variante IR-101 da família ResNet adaptada para faces, é treinado em WebFace4M com uma perda de margem adaptativa que ajusta o peso de exemplos difíceis em função da qualidade da imagem. O resultado é um extrator de 512 dimensões em que distâncias entre embeddings refletem similaridade de identidade. No Modelo 12 deste trabalho, o AdaFace IR-101 entra como backbone facial especializado, com descongelamento parcial do último estágio, e seus embeddings alimentam a estrutura regional descrita no Capítulo~\ref{cap:Metodologia}.

\section{Atenção cruzada entre faces}

Atenção cruzada, ou cross-attention, é uma variante da self-attention em que as consultas vêm de uma sequência e as chaves e valores vêm de outra. No modelo de verificação de parentesco, isso significa que os tokens de patch de uma face geram as consultas e os tokens da outra face geram chaves e valores. Formalmente, dado $X_a$ e $X_b$ os tokens das duas faces, a atenção cruzada de $a$ sobre $b$ é

\begin{equation}
\text{CrossAttn}(X_a, X_b) = \mathrm{softmax}\!\left( \frac{X_a W_Q (X_b W_K)^\top}{\sqrt{d_k}} \right) X_b W_V,
\label{eq:cross_attention}
\end{equation}

onde $W_Q$, $W_K$ e $W_V$ são matrizes aprendíveis. A versão bidirecional alterna o papel das duas faces ao longo de camadas sucessivas, e o efeito prático é que o modelo aprende quais regiões da face $a$ comparar com quais regiões da face $b$, em vez de operar só sobre embeddings globais. Este é o módulo central do FaCoR \cite{su2023facornet} e do Modelo 02 deste trabalho.

Uma variante explorada por este trabalho reduz drasticamente o número de tokens envolvidos na atenção cruzada ao trocar patches densos por regiões anatômicas fixas. Em vez de 196 tokens espaciais como no ViT base, define-se um pequeno conjunto de regiões da face (por exemplo, global, olhos, nariz, boca e mandíbula) e cada região vira um único token. A atenção cruzada entre regiões opera sobre uma matriz pequena, da ordem de cinco por cinco no caso anatômico, e fornece um prior estrutural ao restringir explicitamente onde o modelo compara as duas faces. A formulação completa, descrita no Capítulo~\ref{cap:Metodologia}, combina os tokens regionais com um mecanismo de gate sigmoide aprendido por região, permitindo que múltiplas regiões sejam salientes simultaneamente. Na revisão realizada para este trabalho, não foi identificada combinação equivalente aplicada à verificação de parentesco facial.

O uso de regiões anatômicas fixas também responde a uma limitação prática dos patches densos em bancos de parentesco. Em imagens alinhadas, olhos, nariz, boca e contorno inferior tendem a ocupar faixas relativamente estáveis, mesmo quando há variação de idade, iluminação e expressão. Transformar essas regiões em poucos tokens reduz a complexidade da comparação e força o modelo a operar sobre subestruturas faciais interpretáveis, em vez de depender de todos os patches da imagem, incluindo fundo, cabelo e áreas pouco informativas. O custo dessa escolha é a dependência maior da qualidade do alinhamento facial: se o detector desloca a face ou falha em pose extrema, as regiões fixas deixam de corresponder às partes anatômicas pretendidas.

Outro princípio usado no Modelo 12 é a simetria da verificação de parentesco. A decisão binária ``A é parente de B'' deveria ser equivalente a ``B é parente de A'', isto é, idealmente $f(A,B) = f(B,A)$. Arquiteturas com concatenação ordenada, atenção cruzada direcional ou cabeças que recebem os embeddings crus de $A$ e $B$ podem violar essa propriedade e aprender atalhos dependentes da posição da imagem no par. A passagem simétrica trata esse problema como regularização de consistência: o mesmo par é avaliado nas duas ordens, e a perda força as duas direções a produzirem decisões compatíveis. Essa estratégia não adiciona novos parâmetros, mas reduz a liberdade do modelo para memorizar padrões específicos da ordem registrada no conjunto de treino.

Uma variante mais recente, a atenção diferencial \cite{ye2024differential},
divide cada cabeça de atenção em duas metades e computa a atenção como diferença entre dois mapas softmax,

\begin{equation}
\text{DiffAttn}(Q, K, V) = \big( \mathrm{softmax}(Q_1 K_1^\top / \sqrt{d}) - \lambda \, \mathrm{softmax}(Q_2 K_2^\top / \sqrt{d}) \big) V,
\label{eq:diff_attn}
\end{equation}

onde $\lambda$ é um escalar aprendível por cabeça. A intuição é que a primeira metade captura o sinal relevante e a segunda captura ruído correlacionado, e a subtração elimina correspondências espúrias. No domínio facial isso é particularmente útil porque patches de fundo, cabelo e variação de iluminação produzem similaridade alta sem significado de parentesco. O Modelo 05 deste trabalho usa atenção diferencial no lugar da self-attention padrão dentro do módulo de atenção cruzada.

\section{Adaptação de baixo posto com LoRA}

Low-Rank Adaptation, ou LoRA \cite{hu2021lora},
é uma técnica de fine-tuning eficiente em parâmetros. Em vez de re-treinar a matriz de pesos $W \in \mathbb{R}^{d \times d}$ de uma camada linear, com $d^2$ parâmetros, LoRA congela $W$ e injeta um caminho aditivo paralelo de baixo posto,

\begin{equation}
y = Wx + \frac{\alpha}{r} \, B A x,
\label{eq:lora}
\end{equation}

onde $A \in \mathbb{R}^{r \times d}$ e $B \in \mathbb{R}^{d \times r}$ são as matrizes treináveis de baixo posto $r$, com $r \ll d$, e $\alpha$ é um fator de escala constante. Apenas $A$ e $B$ recebem gradiente, o que reduz drasticamente o número de parâmetros treináveis e o consumo de memória. No Modelo 05 deste trabalho, LoRA com $r=8$ e $\alpha=16$ é injetado nas projeções qkv e proj de cada bloco do DINOv2, e adiciona cerca de 2 a 3 milhões de parâmetros treináveis sobre os 86 milhões do backbone congelado. A restrição de baixo posto também introduz regularização implícita, o que ajuda a evitar overfitting quando os dados de fine-tuning são limitados.

\section{Fusão de representações multibackbone}

Modelos com mais de um backbone, como o híbrido ConvNeXt mais ViT do Modelo 03, precisam combinar as representações de cada ramo em um único embedding antes da projeção final. As estratégias de fusão variam em expressividade. A mais simples é a concatenação direta dos vetores de cada ramo, seguida de uma projeção linear. O gating multiplicativo aprende um vetor de pesos por dimensão que pondera as contribuições de cada ramo. A fusão bilinear computa o produto externo entre os dois vetores e captura interações em todas as combinações de atributos. A fusão por atenção trata cada vetor como um token e aprende pesos de mistura por meio de um mecanismo de atenção sobre os dois ramos. Em geral, fusões mais expressivas têm mais parâmetros e mais risco de sobreajuste, e a concatenação simples costuma ser competitiva quando os dados de treino são limitados, padrão observado em diversos estudos comparativos de fusão multimodal e replicado no presente trabalho no Modelo 03.

\section{Modelos aumentados por recuperação}

Modelos aumentados por recuperação combinam pesos aprendíveis com uma memória não-paramétrica externa. A ideia é que parte do conhecimento necessário para a decisão não precisa estar codificada nos parâmetros do modelo e pode ser recuperada de uma base de exemplos no momento da inferência. Esse paradigma é dominante em modelagem de linguagem, em sistemas conhecidos como RAG ou kNN-LM, mas ainda é raro em visão e mais raro ainda em verificação de parentesco facial.

A formulação básica para verificação de parentesco é direta. Dado um par de consulta $(x_a, x_b)$, calcula-se uma assinatura $s_q$ que resume o par no espaço de embeddings. Numa galeria pré-computada $\mathcal{G} = \{(s_i, r_i)\}$, em que $s_i$ é a assinatura de um par positivo do conjunto de treino e $r_i$ é o tipo da relação, recupera-se os $K$ pares mais similares a $s_q$ por cosseno. Os $K$ pares recuperados viram tokens de contexto para um módulo de atenção cruzada que olha da consulta para os suportes, e a saída desse módulo alimenta a cabeça binária e uma cabeça auxiliar de classificação da relação. O Modelo 06 deste trabalho segue exatamente essa receita, com galeria de cerca de 33 mil pares positivos da partição de treino do FIW.

\section{Funções de perda}

A função de perda mais simples para verificação binária é a entropia cruzada binária, BCE, aplicada sobre o logit produzido pela cabeça final do modelo. BCE funciona como critério de classificação direta, mas não impõe explicitamente uma restrição geométrica sobre o espaço de embeddings, e por isso é usada em conjunto com uma família de perdas de metric learning quando o objetivo é manter pares positivos próximos e pares negativos distantes no embedding.

A Contrastive Loss clássica, proposta por \citeonline{hadsell2006contrastive}, trabalha sobre pares de imagens e aproxima positivos enquanto afasta negativos até uma margem $m$. Em sua versão com distância cosseno,

\begin{equation}
\mathcal{L}_{\text{con}} = y \cdot (1 - \cos(z_a, z_b)) + (1 - y) \cdot \max\!\big(0, \cos(z_a, z_b) - m\big),
\label{eq:contrastive}
\end{equation}

em que $y$ é o rótulo do par, com 1 para parente e 0 para não-parente, $z_a$ e $z_b$ são os embeddings $L^2$-normalizados, e $m$ é a margem.

Em vez de pares, a Triplet Loss \citeonline{schroff2015facenet} opera sobre trios formados por uma âncora $z_a$, um positivo $z_p$ que pertence à mesma classe da âncora, e um negativo $z_n$ que pertence a classe diferente. A perda força a distância entre âncora e positivo a ser menor que a distância entre âncora e negativo por uma margem mínima.

A Supervised Contrastive Loss \cite{khosla2020supervised} generaliza a Contrastive Loss para múltiplos positivos por amostra no batch e estende o efeito de regularização do contraste para classificação supervisionada. A formulação canônica dos autores usa similaridade cosseno escalada por temperatura e soma sobre todos os positivos no batch, sem margem explícita. A versão implementada nas execuções referenciadas como Supervised Contrastive no Capítulo~\ref{cap:Metodologia} é uma adaptação operacional do conceito, com distância cosseno entre pares e margem $m = 0{,}3$, mais próxima da Contrastive Loss clássica do que da formulação canônica de Khosla. A escolha foi guiada por estabilidade empírica em FIW e por simplicidade de implementação no pipeline binário dos modelos avaliados.

A escolha entre essas perdas afeta diretamente a geometria do espaço de embeddings aprendido. Perdas com margem produzem clusters mais compactos e mais separáveis, o que é útil tanto para verificação quanto para classificação de relação.

\section{Síntese de idade}

A diferença de idade entre pais e filhos é um confundidor conhecido em verificação de parentesco. Quando o pai aparece com 50 anos e o filho com 20, a similaridade facial é parcialmente mascarada por mudanças associadas ao envelhecimento. Métodos de síntese de idade tentam neutralizar esse efeito, gerando variantes etárias da face de cada indivíduo antes da comparação.

O SAM, sigla em inglês para Style-based Age Manipulation \cite{alaluf2022sam}, é uma implementação dessa ideia baseada em StyleGAN. O método aprende um mapeamento entre o espaço latente do gerador e a idade alvo, e produz, a partir de uma única imagem, variantes do mesmo rosto em diferentes idades. Neste trabalho, o esquema de invariância geracional baseado em três variantes etárias por face (criança, jovem e idosa) foi incorporado apenas ao Modelo 01. Os Modelos 02, 05, 06 e 12 não usam síntese de idade, de modo que a comparação principal não mistura ganhos de arquitetura com ganhos de manipulação etária. A descrição completa do esquema e da sua aplicação restrita está no Capítulo~\ref{cap:Metodologia}.

\section{Modelos multimodais visão-linguagem}

Modelos multimodais visão-linguagem, ou VLMs, combinam um encoder visual com um modelo de linguagem grande. CLIP \cite{radford2021learning} foi um dos primeiros nesse formato e treina o encoder visual e o encoder textual conjuntamente para maximizar a similaridade entre pares imagem-texto que de fato co-ocorrem. Modelos mais recentes, como GPT-4V e Claude Sonnet, ampliam o escopo para diálogo multimodal arbitrário, em que o usuário envia uma imagem junto com uma instrução textual e o modelo responde em linguagem natural.

Para tarefas como verificação de parentesco, a capacidade interessante dos VLMs é o regime zero-shot. O usuário envia a imagem composta de duas faces e um prompt que descreve a tarefa, e o modelo responde sem treinamento específico no domínio. As linhas de base do Capítulo~\ref{cap:Metodologia} modificam esse comportamento sem fine-tuning por três mecanismos combinados. No few-shot in-context learning, o prompt inclui alguns exemplos com a resposta correta para orientar o modelo na tarefa específica. Em torno dele, um prompt estruturado obriga o modelo a produzir saídas intermediárias antes da resposta final, como estimar diferença geracional aparente ou gênero, o que aproxima a inferência de uma cadeia de raciocínio explícita. Sobre essas saídas auxiliares, uma calibração leve treinada em validação ajusta a resposta final.

\section{Métricas de avaliação}

Como métrica principal deste trabalho adota-se a área sob a curva ROC, abreviada ROC AUC, que mede a probabilidade de o modelo atribuir um escore mais alto a um par positivo aleatório do que a um par negativo aleatório. O AUC tem a vantagem de ser invariante ao limiar de decisão, o que permite comparar modelos sem depender de uma escolha de $\tau$; valores próximos de 1 indicam separação perfeita e valores próximos de 0,5 indicam classificador aleatório.

As demais métricas dependem do limiar escolhido. Acurácia mede a fração de classificações corretas no conjunto de teste, enquanto precisão e revocação medem, respectivamente, a fração de pares preditos como positivos que de fato são positivos e a fração de pares positivos que foram corretamente identificados, e o F1 equilibra os dois numa única métrica como média harmônica entre precisão e revocação. Essas métricas são reportadas com o limiar $\tau^*$ selecionado em validação, conforme o protocolo descrito no Capítulo~\ref{cap:Metodologia}.

Para aplicações em que falsos positivos têm custo alto, como forense, usa-se TAR@FAR, ou taxa de aceitação verdadeira em uma taxa fixa de falsa aceitação. Este trabalho reporta TAR@FAR em $0,1$, $0,01$ e $0,001$, que correspondem respectivamente a um regime tolerante, a um regime estrito e a um regime extremamente estrito de falsa aceitação. Average Precision, por sua vez, integra precisão sobre a curva precisão-revocação e captura o desempenho médio em todos os regimes de limiar simultaneamente.

Para as linhas de base com VLMs, a tarefa é multiclasse com onze classes, e as métricas adequadas são acurácia global, macro-precision, macro-recall e macro-F1. A versão macro promedia o resultado por classe, sem ponderar pelo número de amostras, e por isso revela bem o comportamento em classes raras como avós e netos. A matriz de confusão complementa o panorama mostrando para quais classes erradas o modelo desvia cada classe correta.
